% answers.tex —— 2012 年·福建卷·文科 · 参考答案与详细解析片段
% 由 tools/generate_projects.py 依来源总账自 2012/2012_数学真题_福建文_试题.md + 2012/2012_数学真题_福建文_解析正文.md + 2012/2012_数学真题_福建文_答案要点.md 生成，勿手改（改源 MD 后重跑）


\begin{answer}
\gksolhead{1}{}

\textbf{题目}：

复数 $(2+\mathrm i)^2$ 等于\mbox{（　　）}

（A）$3+4\mathrm i$；（B）$5+4\mathrm i$；（C）$3+2\mathrm i$；（D）$5+2\mathrm i$。

\Needspace{4\baselineskip}
\textbf{解析}

\Needspace{5\baselineskip}
按复数乘法展开，并用 $\mathrm i^2=-1$：
\nopagebreak[4]
\[\begin{aligned}
(2+\mathrm i)^2&=2^2+2\cdot2\cdot\mathrm i+\mathrm i^2\\
&=4+4\mathrm i-1\\
&=3+4\mathrm i.
\end{aligned}\]
实部为 $3$、虚部为 $4$。

选 A。

\gkanswer{答案：A。$(2+\mathrm i)^2=4+4\mathrm i+\mathrm i^2=3+4\mathrm i$。}

\end{answer}
\begin{answer}
\gksolhead{2}{}

\textbf{题目}：

已知集合 $M=\{1,2,3,4\}$，$N=\{-2,2\}$，下列结论成立的是\mbox{（　　）}

（A）$N\subseteq M$；（B）$M\cup N=M$；（C）$M\cap N=N$；（D）$M\cap N=\{2\}$。

\Needspace{4\baselineskip}
\textbf{解析}

逐项核对。因为 $-2\in N$ 而 $-2\notin M$，所以 $N$ 不是 $M$ 的子集，（A）不成立。

\Needspace{5\baselineskip}
并集为
\nopagebreak[4]
\[M\cup N=\{-2,1,2,3,4\}\neq M,\]
（B）不成立。

\Needspace{5\baselineskip}
交集为
\nopagebreak[4]
\[M\cap N=\{2\},\]
它不等于 $N=\{-2,2\}$，（C）不成立；而（D）正是 $M\cap N=\{2\}$。

选 D。

\gkanswer{答案：D。$-2\in N$ 但 $-2\notin M$，故 $N\not\subseteq M$、$M\cup N\neq M$、$M\cap N\neq N$；而 $M\cap N=\{2\}$ 成立。}

\end{answer}
\begin{answer}
\gksolhead{3}{}

\textbf{题目}：

已知向量 $\boldsymbol a=(x-1,2)$，$\boldsymbol b=(2,1)$，则 $\boldsymbol a\perp\boldsymbol b$ 的充要条件是\mbox{（　　）}

（A）$x=-\dfrac12$；（B）$x=-1$；（C）$x=5$；（D）$x=0$。

\Needspace{4\baselineskip}
\textbf{解析}

\Needspace{5\baselineskip}
两向量垂直的充要条件是数量积为零：
\nopagebreak[4]
\[\boldsymbol a\cdot\boldsymbol b=(x-1)\cdot2+2\cdot1=2x-2+2=2x.\]
\Needspace{5\baselineskip}
令 $2x=0$，得
\nopagebreak[4]
\[x=0.\]
反过来，$x=0$ 时 $\boldsymbol a=(-1,2)$，$\boldsymbol a\cdot\boldsymbol b=-2+2=0$，确为垂直，故 $x=0$ 是充要条件。

选 D。

\gkanswer{答案：D。$\boldsymbol a\perp\boldsymbol b\iff\boldsymbol a\cdot\boldsymbol b=2(x-1)+2=2x=0$，即 $x=0$。}

\end{answer}
\begin{answer}
\gksolhead{4}{}

\textbf{题目}：

一个几何体的三视图形状都相同，大小均相等，那么这个几何体不可以是\mbox{（　　）}

（A）球；（B）三棱锥；（C）正方体；（D）圆柱。

\Needspace{4\baselineskip}
\textbf{解析}

三视图是全等的正投影，要求正视图、侧视图、俯视图形状相同且大小相等。

\begin{itemize}
\item 球：三个视图都是半径相同的圆，可以；
\item 正方体：三个视图都是全等的正方形，可以；
\item 三棱锥：取顶点为 $(1,1,1)$，$(1,-1,-1)$，$(-1,1,-1)$，$(-1,-1,1)$ 的四面体，它向三个坐标面投影都得到同一个“正方形加两条对角线”的图形，三个视图全同，可以；
\item 圆柱：正视图与侧视图是同样大小的矩形（宽为底面直径、高为柱高），俯视图是圆，矩形不可能与圆形状相同、大小相等，不可以。
\end{itemize}

选 D。

\gkanswer{答案：D。圆柱的正视图与侧视图是矩形、俯视图是圆，三视图不全相同；球、三棱锥、正方体都可做到三视图全同。}

\end{answer}
\begin{answer}
\gksolhead{5}{}

\textbf{题目}：

已知双曲线 $\dfrac{x^2}{a^2}-\dfrac{y^2}{5}=1$ 的右焦点为 $(3,0)$，则该双曲线的离心率等于\mbox{（　　）}

（A）$\dfrac{3\sqrt{14}}{14}$；（B）$\dfrac{3\sqrt2}{4}$；（C）$\dfrac32$；（D）$\dfrac43$。

\Needspace{4\baselineskip}
\textbf{解析}

\Needspace{5\baselineskip}
双曲线 $\dfrac{x^2}{a^2}-\dfrac{y^2}{5}=1$ 中 $b^2=5$。右焦点为 $(3,0)$，故半焦距
\nopagebreak[4]
\[c=3.\]
\Needspace{5\baselineskip}
由 $c^2=a^2+b^2$ 得
\nopagebreak[4]
\[a^2=c^2-b^2=9-5=4,\]
\Needspace{5\baselineskip}
所以 $a=2$。离心率
\nopagebreak[4]
\[e=\frac ca=\frac32.\]

选 C。

\gkanswer{答案：C。由 $c=3$ 及 $c^2=a^2+5$ 得 $a^2=4$，$a=2$，故 $e=\dfrac ca=\dfrac32$。}

\end{answer}
\begin{answer}
\gksolhead{6}{}

\textbf{题目}：

阅读如图所示的程序框图，运行相应的程序，输出的 $s$ 值等于\mbox{（　　）}
\nopagebreak[4]
\begin{center}% ============================================================
% f2012_fujian_arts_q06_flowchart.tex
% 2012 福建卷（文科）第 6 题 程序框图（TikZ 矢量重绘）
% 开始 → k=1,s=1 → 判断 k<4
%   是：s=2s-k → k=k+1 → 回到判断
%   否：输出 s → 结束
% 来源 PDF：普通高考/2012/2012福建文.pdf，印面 731 第 1 栏。
% ============================================================
\begin{tikzpicture}[
  x=1cm, y=1cm, line width=0.6pt, font=\footnotesize, >=Stealth,
  every node/.style={inner sep=3pt},
  proc/.style={rectangle, draw, minimum width=2.7cm, minimum height=0.6cm},
  io/.style={trapezium, draw, trapezium left angle=70, trapezium right angle=110,
             minimum width=2.5cm, minimum height=0.62cm, trapezium stretches=true},
  dec/.style={diamond, draw, aspect=2.3, minimum width=2.2cm, minimum height=0.95cm},
  term/.style={rectangle, rounded corners=8pt, draw, minimum width=1.7cm, minimum height=0.56cm}
]

  \node[term] (start) at (0,0) {开始};
  \node[proc] (init)  at (0,-1.25) {$k=1$，$s=1$};
  \node[dec]  (judge) at (0,-2.95) {$k<4$?};

  \node[proc] (sbox)  at (3.05,-2.35) {$s=2s-k$};
  \node[proc] (kbox)  at (3.05,-1.45) {$k=k+1$};

  \node[io]   (out)   at (0,-4.55) {输出 $s$};
  \node[term] (end)   at (0,-5.65) {结束};

  \draw[->] (start) -- (init);
  \draw[->] (init) -- (judge);
  \draw[->] (judge.east) -- (3.05,-2.95) -- (sbox.south);
  \draw[->] (sbox.north) -- (kbox.south);
  \draw[->] (kbox.north) -- (3.05,-0.92) -- (0,-0.92);
  \draw[->] (judge.south) -- (out.north);
  \draw[->] (out) -- (end);

  \node at (1.62,-2.74) {是};
  \node at (0.38,-3.72) {否};

\end{tikzpicture}
\end{center}
\nopagebreak[4]
（A）$-3$；（B）$-10$；（C）$0$；（D）$-2$。

\Needspace{4\baselineskip}
\textbf{解析}

\Needspace{5\baselineskip}
初值 $k=1$，$s=1$。每次循环先算 $s=2s-k$，再算 $k=k+1$，然后判断 $k<4$：

\nopagebreak[4]
\[\begin{aligned}
k=1&:\quad s=2\times1-1=1,\quad k=2,\\
k=2&:\quad s=2\times1-2=0,\quad k=3,\\
k=3&:\quad s=2\times0-3=-3,\quad k=4.
\end{aligned}\]

\Needspace{5\baselineskip}
此时 $k=4$，条件 $k<4$ 不成立，退出循环，输出
\nopagebreak[4]
\[s=-3.\]

选 A。

\gkanswer{答案：A。循环体依次算得 $(s,k)=(1,2),(0,3),(-3,4)$，此时 $k<4$ 不成立，输出 $s=-3$。}

\end{answer}
\begin{answer}
\gksolhead{7}{}

\textbf{题目}：

直线 $x+\sqrt3y-2=0$ 与圆 $x^2+y^2=4$ 相交于 $A$，$B$ 两点，则弦 $AB$ 的长度等于\mbox{（　　）}

（A）$2\sqrt5$；（B）$2\sqrt3$；（C）$\sqrt3$；（D）$1$。

\Needspace{4\baselineskip}
\textbf{解析}

\Needspace{5\baselineskip}
圆 $x^2+y^2=4$ 的圆心为 $O(0,0)$，半径 $r=2$。圆心到直线 $x+\sqrt3y-2=0$ 的距离
\nopagebreak[4]
\[d=\frac{|0+0-2|}{\sqrt{1^2+(\sqrt3)^2}}=\frac{2}{2}=1.\]
\Needspace{5\baselineskip}
由垂径定理，弦长
\nopagebreak[4]
\[|AB|=2\sqrt{r^2-d^2}=2\sqrt{4-1}=2\sqrt3.\]

选 B。

\gkanswer{答案：B。圆心到直线距离 $d=\dfrac{|{-2}|}{\sqrt{1+3}}=1$，故 $|AB|=2\sqrt{4-1}=2\sqrt3$。}

\end{answer}
\begin{answer}
\gksolhead{8}{}

\textbf{题目}：

函数 $f(x)=\sin\left(x-\dfrac\pi4\right)$ 的图象的一条对称轴是\mbox{（　　）}

（A）$x=\dfrac\pi4$；（B）$x=\dfrac\pi2$；（C）$x=-\dfrac\pi4$；（D）$x=-\dfrac\pi2$。

\Needspace{4\baselineskip}
\textbf{解析}

\Needspace{5\baselineskip}
$y=\sin u$ 的对称轴为 $u=\dfrac\pi2+k\pi$（$k\in\mathbf Z$）。令
\nopagebreak[4]
\[x-\frac\pi4=\frac\pi2+k\pi,\]
\Needspace{5\baselineskip}
解得
\nopagebreak[4]
\[x=\frac{3\pi}{4}+k\pi.\]
在给出的选项中，取 $k=-1$ 得 $x=-\dfrac\pi4$，其余三项都不在解集中。

选 C。

\gkanswer{答案：C。对称轴满足 $x-\dfrac\pi4=\dfrac\pi2+k\pi$，即 $x=\dfrac{3\pi}4+k\pi$；取 $k=-1$ 得 $x=-\dfrac\pi4$。}

\end{answer}
\begin{answer}
\gksolhead{9}{}

\textbf{题目}：

\Needspace{5\baselineskip}
设
\nopagebreak[4]
\[f(x)=\begin{cases}1,&x>0,\\0,&x=0,\\-1,&x<0,\end{cases}\qquad g(x)=\begin{cases}1,&x\text{ \text{为有理数}},\\0,&x\text{ \text{为无理数}},\end{cases}\]
则 $f(g(\pi))$ 的值为\mbox{（　　）}

（A）$1$；（B）$0$；（C）$-1$；（D）$\pi$。

\Needspace{4\baselineskip}
\textbf{解析}

\Needspace{5\baselineskip}
$\pi$ 是无理数，所以按 $g$ 的定义
\nopagebreak[4]
\[g(\pi)=0.\]
\Needspace{5\baselineskip}
再按 $f$ 的定义，$x=0$ 时 $f(x)=0$，故
\nopagebreak[4]
\[f(g(\pi))=f(0)=0.\]

选 B。

\gkanswer{答案：B。$\pi$ 是无理数，故 $g(\pi)=0$，于是 $f(g(\pi))=f(0)=0$。}

\end{answer}
\begin{answer}
\gksolhead{10}{}

\textbf{题目}：

\Needspace{5\baselineskip}
若直线 $y=2x$ 上存在点 $(x,y)$ 满足约束条件
\nopagebreak[4]
\[\begin{cases}x+y-3\leqslant0,\\x-2y-3\leqslant0,\\x\geqslant m,\end{cases}\]
则实数 $m$ 的最大值为\mbox{（　　）}

（A）$-1$；（B）$1$；（C）$\dfrac32$；（D）$2$。

\Needspace{4\baselineskip}
\textbf{解析}

\Needspace{5\baselineskip}
直线 $y=2x$ 上的点可写成 $(x,2x)$。代入前两个约束：
\nopagebreak[4]
\[\begin{cases}
x+2x-3\leqslant0,\\
x-4x-3\leqslant0,
\end{cases}
\quad\Longleftrightarrow\quad
\begin{cases}
x\leqslant1,\\
-3x\leqslant3,
\end{cases}
\quad\Longleftrightarrow\quad
-1\leqslant x\leqslant1.\]
\Needspace{5\baselineskip}
再加上 $x\geqslant m$，直线 $y=2x$ 上存在满足全部条件的点，当且仅当区间 $[-1,1]$ 与 $[m,+\infty)$ 有公共点，即
\nopagebreak[4]
\[m\leqslant1.\]
所以 $m$ 的最大值为 $1$。

选 B。

\gkanswer{答案：B。把 $y=2x$ 代入前两个不等式得 $-1\leqslant x\leqslant1$，故存在这样的点当且仅当 $m\leqslant1$，$m$ 的最大值为 $1$。}

\end{answer}
\begin{answer}
\gksolhead{11}{}

\textbf{题目}：

数列 $\{a_n\}$ 的通项公式 $a_n=n\cos\dfrac{n\pi}{2}$，其前 $n$ 项和为 $S_n$，则 $S_{2012}$ 等于\mbox{（　　）}

（A）$1006$；（B）$2012$；（C）$503$；（D）$0$。

\Needspace{4\baselineskip}
\textbf{解析}

\Needspace{5\baselineskip}
按 $n$ 被 $4$ 除的余数分类。当 $n=4k-3,4k-1$ 时 $\cos\dfrac{n\pi}{2}=0$；当 $n=4k-2$ 时 $\cos\dfrac{n\pi}{2}=-1$；当 $n=4k$ 时 $\cos\dfrac{n\pi}{2}=1$。于是每一组连续四项
\nopagebreak[4]
\[a_{4k-3}+a_{4k-2}+a_{4k-1}+a_{4k}=0-(4k-2)+0+4k=2.\]
\Needspace{5\baselineskip}
因为
\nopagebreak[4]
\[2012=4\times503,\]
\Needspace{5\baselineskip}
共有 $503$ 个完整的四项组，所以
\nopagebreak[4]
\[S_{2012}=503\times2=1006.\]

选 A。

\gkanswer{答案：A。每连续四项之和为 $0-(4k-2)+0+4k=2$，而 $2012=4\times503$，故 $S_{2012}=2\times503=1006$。}

\end{answer}
\begin{answer}
\gksolhead{12}{}

\textbf{题目}：

已知 $f(x)=x^3-6x^2+9x-abc$，$a<b<c$，且 $f(a)=f(b)=f(c)=0$。现给出如下结论：① $f(0)f(1)>0$；② $f(0)f(1)<0$；③ $f(0)f(3)>0$；④ $f(0)f(3)<0$。其中正确结论的序号是\mbox{（　　）}

（A）①③；（B）①④；（C）②③；（D）②④。

\Needspace{4\baselineskip}
\textbf{解析}

\Needspace{5\baselineskip}
$f$ 是首项系数为 $1$ 的三次多项式，且以 $a<b<c$ 为根，故
\nopagebreak[4]
\[f(x)=(x-a)(x-b)(x-c).\]

\Needspace{5\baselineskip}
求导：
\nopagebreak[4]
\[f'(x)=3x^2-12x+9=3(x-1)(x-3).\]
\Needspace{5\baselineskip}
所以 $f$ 在 $(-\infty,1)$ 上递增，在 $(1,3)$ 上递减，在 $(3,+\infty)$ 上递增，$x=1$ 为极大值点，$x=3$ 为极小值点。三个根互不相同，等价于极大值 $f(1)>0$、极小值 $f(3)<0$。由此并结合 $x\to-\infty$ 时 $f\to-\infty$、$x\to+\infty$ 时 $f\to+\infty$，根的分布为
\nopagebreak[4]
\[a<1<b<3<c.\]

\Needspace{5\baselineskip}
比较 $x=0$ 与 $x=3$ 的函数值。由 $f$ 的展开式，
\nopagebreak[4]
\[f(0)=-abc,\qquad f(3)=27-54+27-abc=-abc,\]
\Needspace{5\baselineskip}
故
\nopagebreak[4]
\[f(0)=f(3).\]
\Needspace{5\baselineskip}
而 $f(3)<0$，所以
\nopagebreak[4]
\[f(0)<0.\]
\Needspace{5\baselineskip}
另一方面 $f(1)>0$。于是
\nopagebreak[4]
\[f(0)f(1)<0,\qquad f(0)f(3)=f(0)^2>0.\]
（$f(0)\neq0$：若 $f(0)=0$，则 $0$ 与 $3$ 都是根，由根与系数关系 $a+b+c=6$、$ab+bc+ca=9$ 可推出另一根也是 $3$，与 $a<b<c$ 矛盾。）

所以②③正确。

选 C。

\gkanswer{答案：C。$f'(x)=3(x-1)(x-3)$，由 $f(1)>0>f(3)$ 得 $a<1<b<3<c$；又 $f(0)=f(3)=-abc<0$，故 $f(0)f(1)<0$、$f(0)f(3)>0$，即②③正确。}

\end{answer}
\begin{answer}
\gksolhead{13}{}

\textbf{题目}：

在 $\triangle ABC$ 中，已知 $\angle BAC=60^\circ$，$\angle ABC=45^\circ$，$BC=\sqrt3$，则 $AC=$\gkblank{}。

\Needspace{4\baselineskip}
\textbf{解析}

\Needspace{5\baselineskip}
在 $\triangle ABC$ 中，$BC$ 是 $\angle A$ 的对边，$AC$ 是 $\angle B$ 的对边。由正弦定理
\nopagebreak[4]
\[\frac{AC}{\sin B}=\frac{BC}{\sin A}.\]
\Needspace{5\baselineskip}
代入 $\angle A=60^\circ$，$\angle B=45^\circ$，$BC=\sqrt3$：
\nopagebreak[4]
\[AC=\frac{BC\sin B}{\sin A}
=\frac{\sqrt3\cdot\frac{\sqrt2}{2}}{\frac{\sqrt3}{2}}
=\sqrt2.\]

\gkanswer{答案：$\sqrt2$。由正弦定理 $AC=\dfrac{BC\sin45^\circ}{\sin60^\circ}=\dfrac{\sqrt3\cdot\frac{\sqrt2}2}{\frac{\sqrt3}2}=\sqrt2$。}

\end{answer}
\begin{answer}
\gksolhead{14}{}

\textbf{题目}：

一支田径队有男女运动员 $98$ 人，其中男运动员有 $56$ 人。按男女比例用分层抽样的方法，从全体运动员中抽出一个容量为 $28$ 的样本，那么应抽取女运动员人数是\gkblank{}。

\Needspace{4\baselineskip}
\textbf{解析}

\Needspace{5\baselineskip}
女运动员人数为
\nopagebreak[4]
\[98-56=42.\]
\Needspace{5\baselineskip}
分层抽样按男女比例分配样本容量，故应抽取的女运动员人数为
\nopagebreak[4]
\[28\times\frac{42}{98}=28\times\frac37=12.\]

\gkanswer{答案：$12$。女运动员 $98-56=42$ 人，按比例应抽 $28\times\dfrac{42}{98}=12$ 人。}

\end{answer}
\begin{answer}
\gksolhead{15}{}

\textbf{题目}：

已知关于 $x$ 的不等式 $x^2-ax+2a>0$ 在 $\mathbf R$ 上恒成立，则实数 $a$ 的取值范围是\gkblank{}。

\Needspace{4\baselineskip}
\textbf{解析}

\Needspace{5\baselineskip}
二次函数 $y=x^2-ax+2a$ 的二次项系数 $1>0$，开口向上。它在 $\mathbf R$ 上恒正，当且仅当判别式小于零：
\nopagebreak[4]
\[\Delta=(-a)^2-4\cdot1\cdot2a=a^2-8a<0.\]
\Needspace{5\baselineskip}
因式分解得
\nopagebreak[4]
\[a(a-8)<0,\]
\Needspace{5\baselineskip}
解得
\nopagebreak[4]
\[0<a<8.\]
即实数 $a$ 的取值范围是 $(0,8)$。

\gkanswer{答案：$(0,8)$。二次项系数为正，恒成立当且仅当判别式 $a^2-8a<0$，解得 $0<a<8$。}

\end{answer}
\begin{answer}
\gksolhead{16}{}

\textbf{题目}：

某地区规划道路建设，考虑道路铺设方案。方案设计图中，点表示城市，两点之间连线表示两城市间可铺设道路，连线上数据表示两城市间铺设道路的费用，要求从任一城市都能到达其余各城市，并且铺设道路的总费用最小。例如：在三个城市道路设计中，若城市间可铺设道路的线路图如图 1，则最优设计方案如图 2，此时铺设道路的最小总费用为 $10$。现给出该地区可铺设道路的线路图如图 3，则铺设道路的最小总费用为\gkblank{}。
\nopagebreak[4]
\begin{center}% ============================================================
% f2012_fujian_arts_q16_fig1.tex
% 2012 福建卷（文科）第 16 题 图 1：三城市线路图
% A、B、C 三城市，AB=4，AC=6，BC=8（数据为道路铺设费用）
% 来源 PDF：普通高考/2012/2012福建文.pdf，印面 731 第 2 栏下部。
% ============================================================
\begin{tikzpicture}[>=Stealth, line width=0.7pt, font=\small, x=1cm, y=1cm]
  \coordinate (A) at (1.45,1.60);
  \coordinate (B) at (0.05,1.05);
  \coordinate (C) at (3.15,0.85);

  \draw (A) -- (B);
  \draw (A) -- (C);
  \draw (B) -- (C);

  \fill (A) circle (1.4pt) node[above] {$A$};
  \fill (B) circle (1.4pt) node[below left] {$B$};
  \fill (C) circle (1.4pt) node[right] {$C$};

  \node at (0.62,1.50) {$4$};
  \node at (2.40,1.42) {$6$};
  \node at (1.55,0.78) {$8$};
  \node at (1.55,-0.55) {图 1};
\end{tikzpicture}
\end{center}
\nopagebreak[4]
\nopagebreak[4]
\begin{center}% ============================================================
% f2012_fujian_arts_q16_fig2.tex
% 2012 福建卷（文科）第 16 题 图 2：图 1 的最优设计方案
% 连接 AB（费用 4）与 AC（费用 6），总费用 10
% 来源 PDF：普通高考/2012/2012福建文.pdf，印面 731 第 2 栏下部。
% ============================================================
\begin{tikzpicture}[>=Stealth, line width=0.7pt, font=\small, x=1cm, y=1cm]
  \coordinate (A) at (1.55,1.65);
  \coordinate (B) at (0.10,0.85);
  \coordinate (C) at (3.40,1.00);

  \draw (A) -- (B);
  \draw (A) -- (C);

  \fill (A) circle (1.4pt) node[above] {$A$};
  \fill (B) circle (1.4pt) node[below left] {$B$};
  \fill (C) circle (1.4pt) node[right] {$C$};

  \node at (0.62,1.40) {$4$};
  \node at (2.55,1.45) {$6$};
  \node at (1.60,-0.35) {图 2};
\end{tikzpicture}
\end{center}
\nopagebreak[4]
\nopagebreak[4]
\begin{center}% ============================================================
% f2012_fujian_arts_q16_fig3.tex
% 2012 福建卷（文科）第 16 题 图 3：七城市道路线路图
% 边权：AB=9，AF=3，AE=2，BF=7，BC=5，FG=1，FE=3，
%       GC=3，GD=2，GE=6，CD=6，DE=6
% 来源 PDF：普通高考/2012/2012福建文.pdf，印面 731 第 2 栏下部。
% ============================================================
\begin{tikzpicture}[>=Stealth, line width=0.7pt, font=\small, x=1cm, y=1cm]
  \coordinate (A) at (0.00,2.55);
  \coordinate (B) at (1.85,3.35);
  \coordinate (C) at (4.35,2.35);
  \coordinate (D) at (3.65,0.55);
  \coordinate (E) at (0.65,0.00);
  \coordinate (F) at (1.55,1.85);
  \coordinate (G) at (2.75,1.55);

  \draw (A) -- (B);
  \draw (A) -- (F);
  \draw (A) -- (E);
  \draw (B) -- (F);
  \draw (B) -- (C);
  \draw (F) -- (G);
  \draw (F) -- (E);
  \draw (G) -- (C);
  \draw (G) -- (D);
  \draw (G) -- (E);
  \draw (C) -- (D);
  \draw (D) -- (E);

  \fill (A) circle (1.3pt) node[left] {$A$};
  \fill (B) circle (1.3pt) node[above] {$B$};
  \fill (C) circle (1.3pt) node[right] {$C$};
  \fill (D) circle (1.3pt) node[right] {$D$};
  \fill (E) circle (1.3pt) node[below left] {$E$};
  \fill (F) circle (1.3pt) node[below left] {$F$};
  \fill (G) circle (1.3pt) node[right] {$G$};

  \node at (0.88,3.16) {$9$};
  \node at (0.55,2.18) {$3$};
  \node at (0.18,1.25) {$2$};
  \node at (1.90,2.55) {$7$};
  \node at (3.30,3.02) {$5$};
  \node at (2.16,1.84) {$1$};
  \node at (0.92,1.02) {$3$};
  \node at (3.50,2.14) {$3$};
  \node at (3.42,1.02) {$2$};
  \node at (1.58,0.64) {$6$};
  \node at (4.18,1.42) {$6$};
  \node at (2.18,-0.02) {$6$};

  \node at (2.20,-0.75) {图 3};
\end{tikzpicture}
\end{center}
\nopagebreak[4]
\Needspace{4\baselineskip}
\textbf{解析}

\Needspace{5\baselineskip}
要求用总费用最小的道路网把七个城市全部连通，且不成圈，这是求最小生成树。把图 3 的边按费用从小到大排列：

\nopagebreak[4]
\[\begin{aligned}
&FG=1,\ AE=2,\ GD=2,\ AF=3,\ FE=3,\ GC=3,\\
&BC=5,\ GE=6,\ CD=6,\ DE=6,\ BF=7,\ AB=9.
\end{aligned}\]

按此顺序逐条选取，若加入后成圈则舍弃：

\begin{itemize}
\item 取 $FG=1$，$AE=2$，$GD=2$，$AF=3$；
\item 取 $FE=3$ 会与 $AE$、$AF$ 构成圈 $A-E-F-A$，舍弃；
\item 取 $GC=3$，连通 $C$；
\item 取 $BC=5$，连通 $B$。
\end{itemize}

\Needspace{5\baselineskip}
此时已选六条边，连通了全部七个城市且无圈，故为最小生成树。总费用
\nopagebreak[4]
\[1+2+2+3+3+5=16.\]

\gkanswer{答案：$16$。按费用从小到大取边，依次取 $FG=1$、$AE=2$、$GD=2$、$AF=3$、$GC=3$、$BC=5$，六条边连通七个城市且不成圈，总费用 $16$。}

\end{answer}
\begin{answer}
\gksolhead{17}{}

\textbf{题目}：

在等差数列 $\{a_n\}$ 和等比数列 $\{b_n\}$ 中，$a_1=b_1=1$，$b_4=8$，$\{a_n\}$ 的前 $10$ 项和 $S_{10}=55$。

（1） 求 $a_n$ 和 $b_n$；

（2） 现分别从 $\{a_n\}$ 和 $\{b_n\}$ 的前 $3$ 项中各随机抽取一项，写出相应的基本事件，并求这两项的值相等的概率。

\Needspace{4\baselineskip}
\textbf{解析}

\textbf{（1）}

\Needspace{5\baselineskip}
设等差数列 $\{a_n\}$ 的公差为 $d$。由 $a_1=1$ 及前 $10$ 项和公式
\nopagebreak[4]
\[S_{10}=10a_1+\frac{10\times9}{2}d=10+45d=55,\]
\Needspace{5\baselineskip}
解得
\nopagebreak[4]
\[d=1.\]
\Needspace{5\baselineskip}
故
\nopagebreak[4]
\[a_n=a_1+(n-1)d=1+(n-1)=n.\]

\Needspace{5\baselineskip}
设等比数列 $\{b_n\}$ 的公比为 $q$。由 $b_1=1$，$b_4=8$，
\nopagebreak[4]
\[q^3=\frac{b_4}{b_1}=8,\]
\Needspace{5\baselineskip}
故 $q=2$，
\nopagebreak[4]
\[b_n=b_1q^{n-1}=2^{n-1}.\]

\textbf{（2）}

\Needspace{5\baselineskip}
$\{a_n\}$ 的前 $3$ 项为 $1,2,3$，$\{b_n\}$ 的前 $3$ 项为 $1,2,4$。分别抽取一项，基本事件为有序对
\nopagebreak[4]
\[(1,1),\ (1,2),\ (1,4),\ (2,1),\ (2,2),\ (2,4),\ (3,1),\ (3,2),\ (3,4),\]
共 $3\times3=9$ 个，且每个基本事件等可能。

\Needspace{5\baselineskip}
两项值相等的事件为
\nopagebreak[4]
\[(1,1),\ (2,2),\]
\Needspace{5\baselineskip}
共 $2$ 个。故所求概率
\nopagebreak[4]
\[P=\frac29.\]

\gkanswer{答案（1）：$a_n=n$，$b_n=2^{n-1}$。由 $S_{10}=5(2+9d)=55$ 得 $d=1$；由 $b_1q^3=8$ 得 $q=2$。}

\gkanswer{答案（2）：$\dfrac29$。基本事件共 $3\times3=9$ 个，两项相等的情形为 $(1,1)$、$(2,2)$，共 $2$ 个。}

\end{answer}
\begin{answer}
\gksolhead{18}{}

\textbf{题目}：

某工厂为了对新研发的一种产品进行合理定价，将该产品按事先拟定的价格进行试销，得到如下数据：

\begin{center}
\begin{tabular}{lllllll}
\toprule
单价 $x$（元） & $8$ & $8.2$ & $8.4$ & $8.6$ & $8.8$ & $9$ \\
\midrule
销量 $y$（件） & $90$ & $84$ & $83$ & $80$ & $75$ & $68$ \\
\bottomrule
\end{tabular}
\end{center}

（1） 求回归直线方程 $\hat y=bx+a$，其中 $b=-20$，$a=\bar y-b\bar x$；

（2） 预计在今后的销售中，销量与单价仍然服从（1）中的关系，且该产品的成本是 $4$ 元/件，为使工厂获得最大利润，该产品的单价应定为多少元？（利润 $=$ 销售收入 $-$ 成本）

\Needspace{4\baselineskip}
\textbf{解析}

\textbf{（1）}

\Needspace{5\baselineskip}
由表中数据，
\nopagebreak[4]
\[\bar x=\frac{8+8.2+8.4+8.6+8.8+9}{6}=\frac{51}{6}=8.5,\]
\nopagebreak[4]
\[\bar y=\frac{90+84+83+80+75+68}{6}=\frac{480}{6}=80.\]
\Needspace{5\baselineskip}
已知 $b=-20$，故
\nopagebreak[4]
\[a=\bar y-b\bar x=80-(-20)\times8.5=80+170=250.\]
\Needspace{5\baselineskip}
回归直线方程为
\nopagebreak[4]
\[\hat y=-20x+250.\]

\textbf{（2）}

\Needspace{5\baselineskip}
单价为 $x$ 元时，销量为 $\hat y=-20x+250$ 件，每件利润 $x-4$ 元，故总利润
\nopagebreak[4]
\[\begin{aligned}
L(x)&=(x-4)(-20x+250)\\
&=-20x^2+330x-1000.
\end{aligned}\]
\Needspace{5\baselineskip}
这是开口向下的二次函数，在
\nopagebreak[4]
\[x=-\frac{330}{2\times(-20)}=\frac{330}{40}=8.25\]
\Needspace{5\baselineskip}
处取得最大值，此时
\nopagebreak[4]
\[L(8.25)=(8.25-4)(-165+250)=4.25\times85=361.25.\]
所以为使工厂获得最大利润，单价应定为 $8.25$ 元。

\gkanswer{答案（1）：$\hat y=-20x+250$。$\bar x=8.5$，$\bar y=80$，故 $a=\bar y-b\bar x=80+170=250$。}

\gkanswer{答案（2）：$8.25$ 元。利润 $L=(x-4)(-20x+250)=-20x^2+330x-1000$，在 $x=\dfrac{330}{40}=8.25$ 处取最大值。}

\end{answer}
\begin{answer}
\gksolhead{19}{}

\textbf{题目}：

如图，在长方体 $ABCD-A_1B_1C_1D_1$ 中，$AB=AD=1$，$AA_1=2$，$M$ 为棱 $DD_1$ 上的一点。

（1） 求三棱锥 $A-MCC_1$ 的体积；

（2） 当 $A_1M+MC$ 取得最小值时，求证：$B_1M\perp$ 平面 $MAC$。
\nopagebreak[4]
\begin{center}% ============================================================
% f2012_fujian_arts_q19_cuboid.tex
% 2012 福建卷（文科）第 19 题 长方体 ABCD-A1B1C1D1
% AB=AD=1，AA1=2，M 为棱 DD1 上一点
% 来源 PDF：普通高考/2012/2012福建文.pdf，印面 732 第 1 栏上部。
% ============================================================
\begin{tikzpicture}[>=Stealth, line width=0.7pt, font=\small, x=1cm, y=1cm]
  \coordinate (B) at (0.00,0.00);
  \coordinate (C) at (1.60,0.00);
  \coordinate (B1) at (0.00,2.55);
  \coordinate (C1) at (1.60,2.55);
  \coordinate (A) at (0.95,0.58);
  \coordinate (D) at (2.55,0.58);
  \coordinate (A1) at (0.95,3.13);
  \coordinate (D1) at (2.55,3.13);
  \coordinate (M) at (2.55,1.55);

  % 可见棱
  \draw (B) -- (C) -- (C1) -- (B1) -- cycle;
  \draw (B1) -- (A1) -- (D1) -- (C1);
  \draw (C) -- (D) -- (D1);

  % 不可见棱
  \draw[dashed] (B) -- (A) -- (D);
  \draw[dashed] (A) -- (A1);

  \fill (A1) circle (1.2pt) node[above left] {$A_1$};
  \fill (B1) circle (1.2pt) node[left] {$B_1$};
  \fill (C1) circle (1.2pt) node[right] {$C_1$};
  \fill (D1) circle (1.2pt) node[right] {$D_1$};
  \fill (A) circle (1.2pt) node[below left] {$A$};
  \fill (B) circle (1.2pt) node[below left] {$B$};
  \fill (C) circle (1.2pt) node[below right] {$C$};
  \fill (D) circle (1.2pt) node[right] {$D$};
  \fill (M) circle (1.2pt) node[right] {$M$};
\end{tikzpicture}
\end{center}
\nopagebreak[4]
\Needspace{4\baselineskip}
\textbf{解析}

\textbf{（1）}

\Needspace{5\baselineskip}
建立空间直角坐标系：以 $A$ 为原点，$AB$、$AD$、$AA_1$ 所在直线分别为 $x$、$y$、$z$ 轴。由 $AB=AD=1$，$AA_1=2$，得
\nopagebreak[4]
\[A(0,0,0),\ B(1,0,0),\ C(1,1,0),\ D(0,1,0),\]
\nopagebreak[4]
\[A_1(0,0,2),\ B_1(1,0,2),\ C_1(1,1,2),\ D_1(0,1,2).\]
设 $M(0,1,t)$，其中 $0\leqslant t\leqslant2$。

\Needspace{5\baselineskip}
$\triangle MCC_1$ 的三个顶点都在平面 $y=1$ 内，其中 $CC_1$ 垂直于 $DC$，$M$ 到直线 $CC_1$ 的距离为 $1$。于是
\nopagebreak[4]
\[S_{\triangle MCC_1}=\frac12\cdot CC_1\cdot1=\frac12\times2\times1=1.\]
\Needspace{5\baselineskip}
点 $A(0,0,0)$ 到平面 $y=1$ 的距离为 $1$。所以
\nopagebreak[4]
\[V_{A-MCC_1}=\frac13\cdot S_{\triangle MCC_1}\cdot1=\frac13.\]

\textbf{（2）}

\Needspace{5\baselineskip}
由（1）的坐标，
\nopagebreak[4]
\[A_1M=\sqrt{(0-0)^2+(1-0)^2+(t-2)^2}=\sqrt{1+(2-t)^2},\]
\nopagebreak[4]
\[MC=\sqrt{(1-0)^2+(1-1)^2+(0-t)^2}=\sqrt{1+t^2}.\]
\Needspace{5\baselineskip}
于是
\nopagebreak[4]
\[A_1M+MC=\sqrt{1+(2-t)^2}+\sqrt{1+t^2}.\]
\Needspace{5\baselineskip}
由闵可夫斯基不等式，
\nopagebreak[4]
\[\sqrt{1+(2-t)^2}+\sqrt{1+t^2}
\geqslant\sqrt{(1+1)^2+\bigl((2-t)+t\bigr)^2}
=2\sqrt2,\]
等号成立当且仅当 $(1,2-t)$ 与 $(1,t)$ 同向，即 $2-t=t$，$t=1$。

\Needspace{5\baselineskip}
所以当 $t=1$，即 $M$ 为 $DD_1$ 的中点时，$A_1M+MC$ 取得最小值 $2\sqrt2$。此时
\nopagebreak[4]
\[M(0,1,1).\]

\Needspace{5\baselineskip}
下面证明 $B_1M\perp$ 平面 $MAC$。由坐标
\nopagebreak[4]
\[\overrightarrow{MB_1}=(1,0,2)-(0,1,1)=(1,-1,1),\]
\nopagebreak[4]
\[\overrightarrow{MA}=(0,0,0)-(0,1,1)=(0,-1,-1),\]
\nopagebreak[4]
\[\overrightarrow{MC}=(1,1,0)-(0,1,1)=(1,0,-1).\]
\Needspace{5\baselineskip}
计算数量积：
\nopagebreak[4]
\[\overrightarrow{MB_1}\cdot\overrightarrow{MA}=1\times0+(-1)\times(-1)+1\times(-1)=0,\]
\nopagebreak[4]
\[\overrightarrow{MB_1}\cdot\overrightarrow{MC}=1\times1+(-1)\times0+1\times(-1)=0.\]
\Needspace{5\baselineskip}
所以
\nopagebreak[4]
\[MB_1\perp MA,\qquad MB_1\perp MC.\]
\Needspace{5\baselineskip}
$MA$ 与 $MC$ 是平面 $MAC$ 内两条相交于 $M$ 的直线，故
\nopagebreak[4]
\[B_1M\perp\text{\text{平面} }MAC.\]

\gkanswer{答案（1）：$\dfrac13$。$\triangle MCC_1$ 的面积为 $\dfrac12\times2\times1=1$，$A$ 到平面 $MCC_1$ 的距离为 $1$，故 $V=\dfrac13\times1\times1=\dfrac13$。}

\gkanswer{答案（2）：见解析版}

\end{answer}
\begin{answer}
\gksolhead{20}{}

\textbf{题目}：

某同学在一次研究性学习中发现，以下五个式子的值都等于同一个常数：

① $\sin^2 13^\circ+\cos^2 17^\circ-\sin 13^\circ\cos 17^\circ$；

② $\sin^2 15^\circ+\cos^2 15^\circ-\sin 15^\circ\cos 15^\circ$；

③ $\sin^2 18^\circ+\cos^2 12^\circ-\sin 18^\circ\cos 12^\circ$；

④ $\sin^2(-18^\circ)+\cos^2 48^\circ-\sin(-18^\circ)\cos 48^\circ$；

⑤ $\sin^2(-25^\circ)+\cos^2 55^\circ-\sin(-25^\circ)\cos 55^\circ$。

（1） 试从上述五个式子中选择一个，求出这个常数；

（2） 根据（1）的计算结果，将该同学的发现推广为三角恒等式，并证明你的结论。

\Needspace{4\baselineskip}
\textbf{解析}

\textbf{（1）}

\Needspace{5\baselineskip}
选择②。由 $\sin^2\theta+\cos^2\theta=1$ 及二倍角公式 $\sin15^\circ\cos15^\circ=\dfrac12\sin30^\circ$，
\nopagebreak[4]
\[\begin{aligned}
&\sin^2 15^\circ+\cos^2 15^\circ-\sin15^\circ\cos15^\circ\\
={}&1-\frac12\sin30^\circ
=1-\frac12\times\frac12
=\frac34.
\end{aligned}\]
所以这个常数为 $\dfrac34$。

\textbf{（2）}

\Needspace{5\baselineskip}
观察五个式子：两个角之和都是 $30^\circ$，即若第一角为 $\alpha$，第二角为 $30^\circ-\alpha$。推广为恒等式
\nopagebreak[4]
\[\sin^2\alpha+\cos^2(30^\circ-\alpha)-\sin\alpha\cos(30^\circ-\alpha)=\frac34.\]

\Needspace{5\baselineskip}
证明如下。记 $\beta=30^\circ-\alpha$，则左边
\nopagebreak[4]
\[E=\sin^2\alpha+\cos^2\beta-\sin\alpha\cos\beta.\]
\Needspace{5\baselineskip}
用降幂公式与积化和差：
\nopagebreak[4]
\[\sin^2\alpha=\frac{1-\cos2\alpha}{2},\qquad
\cos^2\beta=\frac{1+\cos2\beta}{2},\]
\nopagebreak[4]
\[\sin\alpha\cos\beta=\frac{\sin(\alpha+\beta)+\sin(\alpha-\beta)}{2}
=\frac{\sin30^\circ+\sin(2\alpha-30^\circ)}{2}.\]
\Needspace{5\baselineskip}
于是
\nopagebreak[4]
\[\begin{aligned}
E&=\frac{1-\cos2\alpha}{2}+\frac{1+\cos2\beta}{2}
-\frac{\sin30^\circ+\sin(2\alpha-30^\circ)}{2}\\
&=1+\frac{\cos2\beta-\cos2\alpha}{2}-\frac14-\frac{\sin(2\alpha-30^\circ)}{2}.
\end{aligned}\]
\Needspace{5\baselineskip}
再由和差化积，
\nopagebreak[4]
\[\cos2\beta-\cos2\alpha
=\cos(60^\circ-2\alpha)-\cos2\alpha
=-2\sin30^\circ\sin(30^\circ-2\alpha)
=\sin(2\alpha-30^\circ).\]
\Needspace{5\baselineskip}
代回得
\nopagebreak[4]
\[E=1-\frac14+\frac{\sin(2\alpha-30^\circ)}{2}-\frac{\sin(2\alpha-30^\circ)}{2}=\frac34.\]
恒等式得证。

\gkanswer{答案（1）：$\dfrac34$。取②：$\sin^215^\circ+\cos^215^\circ-\sin15^\circ\cos15^\circ=1-\dfrac12\sin30^\circ=\dfrac34$。}

\gkanswer{答案（2）：见解析版}

\end{answer}
\begin{answer}
\gksolhead{21}{}

\textbf{题目}：

如图，等边三角形 $OAB$ 的边长为 $8\sqrt3$，且其三个顶点均在抛物线 $E:x^2=2py$（$p>0$）上。

（1） 求抛物线 $E$ 的方程；

（2） 设动直线 $l$ 与抛物线 $E$ 相切于点 $P$，与直线 $y=-1$ 相交于点 $Q$。证明：以 $PQ$ 为直径的圆恒过 $y$ 轴上某定点。
\nopagebreak[4]
\begin{center}% ============================================================
% f2012_fujian_arts_q21_parabola.tex
% 2012 福建卷（文科）第 21 题 抛物线 E: x^2=2py (p>0)
% 等边三角形 OAB 的三个顶点均在抛物线上，A、B 关于 y 轴对称
% 来源 PDF：普通高考/2012/2012福建文.pdf，印面 732 第 1 栏中部。
% ============================================================
\begin{tikzpicture}[>=Stealth, line width=0.7pt, font=\small, x=1cm, y=1cm]
  \draw[->] (-3.0,0) -- (3.0,0) node[right] {$x$};
  \draw[->] (0,-1.0) -- (0,3.0) node[above] {$y$};

  \draw[thick, domain=-2.75:2.75, samples=100] plot (\x,{\x*\x/4});

  \coordinate (A) at (-2.40,1.44);
  \coordinate (B) at (2.40,1.44);
  \draw (A) -- (B);

  \fill (A) circle (1.2pt) node[left] {$A$};
  \fill (B) circle (1.2pt) node[right] {$B$};
  \fill (0,0) circle (1.2pt) node[below] {$O$};
\end{tikzpicture}
\end{center}
\nopagebreak[4]
\Needspace{4\baselineskip}
\textbf{解析}

\textbf{（1）}

\Needspace{5\baselineskip}
抛物线 $E:x^2=2py$（$p>0$）与 $y$ 轴只交于原点，而 $O$ 是三角形顶点且在 $E$ 上，故
\nopagebreak[4]
\[O(0,0).\]

\Needspace{5\baselineskip}
设 $A(x_1,y_1)$，$B(x_2,y_2)$。由 $|OA|=|OB|=8\sqrt3$ 及 $x_i^2=2py_i$，
\nopagebreak[4]
\[2py_i+y_i^2=x_i^2+y_i^2=192\qquad(i=1,2).\]
关于 $y$ 的方程 $y^2+2py-192=0$ 只有一个正根，故 $y_1=y_2=:h$，从而 $x_1^2=x_2^2$。又 $A\neq B$，所以 $x_1=-x_2$，即 $A$、$B$ 关于 $y$ 轴对称。

\Needspace{5\baselineskip}
由 $|AB|=8\sqrt3$ 得 $|x_1|=4\sqrt3$，于是
\nopagebreak[4]
\[48=2ph.\]
\Needspace{5\baselineskip}
再由 $|OA|^2=48+h^2=192$ 得
\nopagebreak[4]
\[h^2=144,\qquad h=12.\]
\Needspace{5\baselineskip}
代入 $48=2p\times12$，得
\nopagebreak[4]
\[p=2.\]
\Needspace{5\baselineskip}
所以抛物线 $E$ 的方程为
\nopagebreak[4]
\[x^2=4y.\]

\textbf{（2）}

\Needspace{5\baselineskip}
设 $P(x_0,y_0)$ 在 $E$ 上，则 $x_0^2=4y_0$，且 $x_0\neq0$（$x_0=0$ 时切线为 $y=0$，与 $y=-1$ 无交点）。对 $x^2=4y$ 在 $P$ 处求导得切线斜率 $\dfrac{x_0}{2}$，切线方程为
\nopagebreak[4]
\[y-y_0=\frac{x_0}{2}(x-x_0),\]
\Needspace{5\baselineskip}
利用 $x_0^2=4y_0$ 化简为
\nopagebreak[4]
\[y=\frac{x_0}{2}x-y_0.\]
\Needspace{5\baselineskip}
令 $y=-1$，得交点 $Q$ 的横坐标
\nopagebreak[4]
\[x_Q=\frac{2(y_0-1)}{x_0}.\]
\Needspace{5\baselineskip}
即
\nopagebreak[4]
\[Q\left(\frac{2(y_0-1)}{x_0},-1\right).\]

\Needspace{5\baselineskip}
取 $y$ 轴上的定点 $F(0,1)$。以 $PQ$ 为直径的圆由满足
\nopagebreak[4]
\[(X-P)\cdot(X-Q)=0\]
\Needspace{5\baselineskip}
的点 $X$ 组成。把 $X=F$ 代入：
\nopagebreak[4]
\[\begin{aligned}
(F-P)\cdot(F-Q)
&=(0-x_0)\left(0-\frac{2(y_0-1)}{x_0}\right)+(1-y_0)\bigl(1-(-1)\bigr)\\
&=2(y_0-1)+2(1-y_0)=0.
\end{aligned}\]
所以 $\angle PFQ=90^\circ$，点 $F(0,1)$ 恒在以 $PQ$ 为直径的圆上。

因此该圆恒过 $y$ 轴上的定点 $(0,1)$。

\gkanswer{答案（1）：$x^2=4y$。由对称性设 $A(-4\sqrt3,h)$，$B(4\sqrt3,h)$，由 $|OA|=8\sqrt3$ 得 $h=12$，代入 $48=2p\cdot12$ 得 $p=2$。}

\gkanswer{答案（2）：定点为 $(0,1)$。设 $P(x_0,y_0)$，切线 $y=\dfrac{x_0}2x-y_0$ 与 $y=-1$ 交于 $Q$，可验证 $(0,1)$ 满足以 $PQ$ 为直径的圆的方程。}

\end{answer}
\begin{answer}
\gksolhead{22}{}

\textbf{题目}：

已知函数 $f(x)=ax\sin x-\dfrac32$（$a\in\mathbf R$），且在 $\left[0,\dfrac\pi2\right]$ 上的最大值为 $\dfrac{\pi-3}{2}$。

（1） 求函数 $f(x)$ 的解析式；

（2） 判断函数 $f(x)$ 在 $(0,\pi)$ 内的零点个数，并加以证明。

\Needspace{4\baselineskip}
\textbf{解析}

\textbf{（1）}

\Needspace{5\baselineskip}
求导：
\nopagebreak[4]
\[f'(x)=a\sin x+ax\cos x=a(\sin x+x\cos x).\]
\Needspace{5\baselineskip}
当 $x\in\left[0,\dfrac\pi2\right]$ 时 $\sin x\geqslant0$，$\cos x\geqslant0$，且
\nopagebreak[4]
\[\sin x+x\cos x\geqslant0,\]
仅在 $x=0$ 处取 $0$（$x=\dfrac\pi2$ 时为 $1>0$）。故 $f$ 在 $\left[0,\dfrac\pi2\right]$ 上单调：

\begin{itemize}
\item 若 $a>0$，$f$ 单调不减，最大值在 $x=\dfrac\pi2$ 处，为
\end{itemize}
\nopagebreak[4]
\[f\left(\frac\pi2\right)=a\cdot\frac\pi2\sin\frac\pi2-\frac32=\frac{a\pi}{2}-\frac32;\]
\begin{itemize}
\item 若 $a<0$，$f$ 单调不增，最大值在 $x=0$ 处，为 $f(0)=-\dfrac32$；
\item 若 $a=0$，$f(x)\equiv-\dfrac32$。
\end{itemize}

\Needspace{5\baselineskip}
已知最大值为 $\dfrac{\pi-3}{2}\neq-\dfrac32$，只能是 $a>0$ 的情形。令
\nopagebreak[4]
\[\frac{a\pi}{2}-\frac32=\frac{\pi-3}{2},\]
\Needspace{5\baselineskip}
即
\nopagebreak[4]
\[\frac{a\pi}{2}=\frac\pi2,\]
\Needspace{5\baselineskip}
解得
\nopagebreak[4]
\[a=1.\]
\Needspace{5\baselineskip}
所以
\nopagebreak[4]
\[f(x)=x\sin x-\frac32.\]

\textbf{（2）}

由（1），$f'(x)=\sin x+x\cos x$。

\Needspace{5\baselineskip}
在 $\left(0,\dfrac\pi2\right)$ 上，$\sin x>0$，$\cos x>0$，故 $f'(x)>0$，$f$ 严格递增。又
\nopagebreak[4]
\[f(0)=-\frac32<0,\qquad
f\left(\frac\pi2\right)=\frac\pi2-\frac32=\frac{\pi-3}{2}>0,\]
由零点存在定理及严格单调性，$f$ 在 $\left(0,\dfrac\pi2\right)$ 内恰有一个零点。

\Needspace{5\baselineskip}
在 $\left(\dfrac\pi2,\pi\right)$ 上，再求导：
\nopagebreak[4]
\[f''(x)=\cos x+\cos x-x\sin x=2\cos x-x\sin x.\]
\Needspace{5\baselineskip}
此时 $\cos x<0$，$\sin x>0$，故 $f''(x)<0$，$f'$ 严格递减。而
\nopagebreak[4]
\[f'\left(\frac\pi2\right)=\sin\frac\pi2+\frac\pi2\cos\frac\pi2=1>0,\]
\nopagebreak[4]
\[f'(\pi)=\sin\pi+\pi\cos\pi=-\pi<0,\]
所以存在唯一的 $x_0\in\left(\dfrac\pi2,\pi\right)$ 使 $f'(x_0)=0$。于是 $f$ 在 $\left(\dfrac\pi2,x_0\right)$ 上递增，在 $(x_0,\pi)$ 上递减。

\Needspace{5\baselineskip}
由于
\nopagebreak[4]
\[f\left(\frac\pi2\right)=\frac{\pi-3}{2}>0,\qquad f(\pi)=\pi\sin\pi-\frac32=-\frac32<0,\]
在 $\left(\dfrac\pi2,x_0\right)$ 上 $f(x)\geqslant f\left(\dfrac\pi2\right)>0$，无零点；在 $(x_0,\pi)$ 上 $f$ 严格递减且由正变负，恰有一个零点。

综上，$f(x)$ 在 $(0,\pi)$ 内共有 $2$ 个零点。

\gkanswer{答案（1）：$f(x)=x\sin x-\dfrac32$。由 $f'(x)=a(\sin x+x\cos x)$ 知 $f$ 在 $\left[0,\dfrac\pi2\right]$ 上单调，最大值只能在端点取得，比较得 $a=1$。}

\gkanswer{答案（2）：$2$ 个。$f$ 在 $\left(0,\dfrac\pi2\right)$ 上单调递增且 $f(0^+)<0<f\left(\dfrac\pi2\right)$，在 $\left(\dfrac\pi2,\pi\right)$ 上先增后减且 $f\left(\dfrac\pi2\right)>0>f(\pi^-)$，各恰有一个零点。}

\end{answer}